%==============================================================================
\chapter{القيود والمتطلبات التقنية}
%==============================================================================

\section{البنية}
بنية \textbf{قابلة للتوسّع، متطوّرة، مرنة، متعدّدة الطبقات، وضعيفة الاقتران}،
تستند إلى \textbf{J2EE} كبنية تحتية رئيسية.

\section{المتطلبات التقنية المفروضة}
\begin{xltabular}{\textwidth}{>{\bfseries}p{4cm} X}
\toprule
\textbf{المتطلب} & \textbf{الوصف} \\
\midrule
\endhead
MVC & فصل النموذج / العرض / المتحكّم. \\
Servlet & معالجة الطلبات من جهة الخادم. \\
JSP & واجهة العرض الأمامية. \\
EJB & مكوّنات الأعمال. \\
JDBC & الوصول إلى المعطيات. \\
التدويل & نصوص الواجهة مُخرَجة في ملف XML، انفتاح على الترجمة دون اتصال دون حدّ للّغات، الفرنسية والإنجليزية والعربية على الأقل في العرض. \\
الإعداد & الضبط عبر ملف نصّي \texttt{ConfigZumm.ini}. \\
الأمان بين المكوّنات & شهادات \textbf{X.509} وبروتوكول \textbf{SSL}. \\
مصادقة المستخدم & \textbf{OAuth Google} عبر خدمة ويب XML وواجهتها البرمجية. \\
الواجهة البرمجية الخارجية & خدمة ويب XML تعرض \texttt{float getZummHoneyActualQuantity(int zummID)} لاسترجاع كمية عسل خلية. \\
الواجهة الأمامية & JavaScript وأطر عمل حرة، مع احترام سهولة الاستعمال المفروضة. \\
\bottomrule
\end{xltabular}

\section{المتطلبات غير الوظيفية}
\begin{itemize}
  \item \textbf{سهولة الاستعمال والودّية والبساطة}، اتّساق الواجهة وجودة التصميم الغرافيكي (أهمية قصوى).
  \item \textbf{القابلية للتطوّر}: انفتاح على إضافة لغات ووحدات ومؤشّرات.
  \item \textbf{قابلية التشغيل البيني}: خدمات ويب قابلة للاستجواب من تطبيقات خارجية.
  \item \textbf{الأمان}: تشفير التبادلات ومصادقة قوية.
  \item \textbf{جودة التصميم}: احترام مبادئ \textbf{SOLID} واعتماد \textbf{أنماط التصميم} لضمان الاقتران الضعيف والقابلية للتطوّر، كما هو مفصّل في الملحق~\ref{chap:solid}.
\end{itemize}

\begin{encadre}
\textbf{مبادئ التصميم المعتمَدة:} تطبّق البنية مبادئ \textbf{SOLID} الخمسة
(المسؤولية الوحيدة، الانفتاح/الانغلاق، استبدال ليسكوف، فصل الواجهات، عكس
التبعية) ومجموعة من \textbf{أنماط التصميم} (Composite، State، Strategy،
Observer، Command، Factory Method، Builder، Adapter، Facade، Singleton، Proxy).
تطبيقها على النظام، المُبرَّر «قبل / بعد»، مفصّل في الملحق~\ref{chap:solid}.
\end{encadre}

\section{أسئلة التصميم الواجب تغطيتها}
يجب أن يجيب التصميم صراحةً عن الأسئلة التالية:
\begin{enumerate}
  \item ما هي \textbf{مُدخلات} (قائمة، قاموس معطيات وأنواع) و\textbf{مُخرجات} (قائمة، قاموس، أنواع، صيَغ واستعلامات حساب) النظام؟
  \item ممّ يتكوّن الحلّ؟
  \item من هم المستخدمون النموذجيون وما حقوق وصولهم إلى الوظائف؟
  \item كيف يُنجز النظام عمليات المستخدم؟
  \item كيف يُحزَم النظام ويُنشَر؟
  \item كيف تتفاعل العناصر وبأيّ ترتيب / تسلسل زمني؟
  \item لماذا وكيف وأين استُعملت كل تقنية؟
\end{enumerate}

\section{تبرير الخيارات التقنية}
يجيب الجدول التالي عن لماذا وكيف وأين كل تقنية مفروضة. \textbf{التنفيذات
الملموسة} لهذا الأساس (خادم التطبيقات، نظام إدارة قواعد المعطيات، مكتبات الواجهة
الأمامية، أسلوب خدمة الويب) و\textbf{المقارنة المُبرَّرة} للبدائل المُقيَّمة
مفصّلة في الملحق~\ref{chap:technologies}.
\begin{xltabular}{\textwidth}{>{\bfseries}p{2.6cm} X p{5.4cm}}
\toprule
\textbf{التقنية} & \textbf{لماذا} & \textbf{أين} \\
\midrule
\endhead
MVC & فصل المسؤوليات وتسهيل الصيانة & التنظيم العام للتطبيق \\
Servlet & معالجة الطلبات وتنسيق المعالجات & طبقة التحكم \\
JSP & إنتاج الصفحات الديناميكية & طبقة العرض \\
EJB & تغليف منطق الأعمال وجعله قابلًا لإعادة الاستعمال & طبقة الأعمال \\
JDBC & الوصول إلى قاعدة المعطيات العلائقية بطريقة معيارية & طبقة الوصول إلى المعطيات \\
XML i18n & إخراج النصوص للترجمة دون اتصال & وحدة التدويل \\
ConfigZumm.ini & ضبط النظام دون إعادة تصريف & وحدة الإعداد \\
OAuth Google & مصادقة المستخدمين دون إدارة كلمات السرّ & الوصول إلى النظام \\
خدمة ويب API XML & عرض المعطيات للتطبيقات الخارجية & واجهة التشغيل البيني \\
X.509 وSSL & تشفير التبادلات ومصادقتها & الاتصالات بين المكوّنات \\
JavaScript & إثراء سهولة استعمال الواجهة الأمامية & طبقة العرض \\
\bottomrule
\end{xltabular}

\section{التحزيم والنشر}
\begin{itemize}
  \item \textbf{التحزيم}: يُحزَم التطبيق في شكل أرشيف Java مؤسّسي (WAR للجزء الويب، EAR إذا جُمعت EJB)، متضمّنًا ملفات التدويل XML وملف \texttt{ConfigZumm.ini}.
  \item \textbf{الخادم}: النشر على خادم تطبيقات J2EE يوفّر حاوية servlet وحاوية EJB.
  \item \textbf{قاعدة المعطيات}: نظام إدارة قواعد معطيات علائقي يُوصَل عبر JDBC.
  \item \textbf{الأمان}: تثبيت شهادات X.509 وتفعيل SSL للتبادلات، وضبط وصول OAuth Google.
  \item \textbf{الإعداد}: تُضبَط العتبات والوحدات والوحدات النشطة في \texttt{ConfigZumm.ini} دون إعادة نشر.
\end{itemize}
يُوضَّح التوزيع المادي بمخطّط النشر (الشكل \ref{fig:deploiement}).
